\clearpage
\thispagestyle{empty}

\chapter*{{\emph\LARGE{Abstract}}}

The development of flexible and interoperable industrial platforms is essential for addressing the growing integration challenges of Industrial Internet of Things (IIoT) systems. This report presents the design and development of the IO-Box, a configurable industrial Input/Output platform intended to simplify the acquisition, processing, control, and exchange of data between sensors, actuators, and supervisory systems.

The developed platform is based on an STM32 microcontroller and integrates multiple industrial communication interfaces, peripheral drivers, and sensor management modules within a modular software architecture. The implementation emphasizes hardware abstraction, software reusability, and scalability, enabling efficient integration of heterogeneous devices while reducing development effort.

The resulting solution provides a versatile and reusable platform for industrial monitoring, automation, and IoT applications. By improving interoperability and simplifying system integration, the IO-Box contributes to the development of more efficient and scalable industrial solutions.

\textbf{Keywords:} Industrial IoT, Embedded Systems, IO-Box, Data Acquisition, Industrial Communication Protocols, Sensors, Automation, Software Architecture.

\clearpage
\chapter*{{\emph\Large{Resume}}}

Le développement de plateformes industrielles flexibles et interopérables est devenu essentiel pour répondre aux défis croissants d’intégration des systèmes de l’Internet Industriel des Objets (IIoT). Ce rapport présente la conception et le développement de l’IO-Box, une plateforme industrielle d’Entrées/Sorties configurable destinée à simplifier l’acquisition, le traitement, le contrôle et l’échange de données entre les capteurs, les actionneurs et les systèmes de supervision.

La plateforme développée repose sur un microcontrôleur STM32 et intègre plusieurs interfaces de communication industrielle, pilotes de périphériques et modules de gestion de capteurs au sein d’une architecture logicielle modulaire. L’implémentation met l’accent sur l’abstraction matérielle, la réutilisabilité logicielle et l’évolutivité, permettant une intégration efficace de dispositifs hétérogènes tout en réduisant les efforts de développement.

La solution obtenue constitue une plateforme polyvalente et réutilisable pour les applications de supervision industrielle, d’automatisation et d’IoT. En améliorant l’interopérabilité et en simplifiant l’intégration des systèmes, l’IO-Box contribue au développement de solutions industrielles plus efficaces et évolutives.

\textbf{Mots-clés :} Internet Industriel des Objets (IIoT), Systèmes Embarqués, STM32, IO-Box, Acquisition de Données, Protocoles de Communication Industrielle, Automatisation.
